\abstract{%
  Cosmochrony is a pre-geometric relational framework that does not replace
  quantum mechanics or general relativity, but operates at a deeper structural
  level from which both emerge as effective descriptions in appropriate regimes.
  Its aim is to expose the common infra-physical foundations underlying these
  theories, to clarify how their conceptual structures are mutually consistent,
  and to show that several phenomena currently treated as independent postulates
  or open problems follow necessarily from a minimal common origin.

  The framework postulates a single static substrate~$\chi$, encoding only
  admissibility constraints, together with a fixed generically non-injective
  projection~$\Pi$ mapping substrate configurations onto observable states.
  No fundamental spacetime, metric, or quantum structure is assumed.
  From this minimal ontology, temporal ordering, spatial geometry, the
  causal structure of spacetime, the origin of quantum indeterminacy, and
  the structure of matter emerge as projective consequences rather than
  independent postulates.
  The speed of light and Planck's constant arise as complementary bounds
  on projectability and spectral resolvability.
  The three-generation structure of Standard Model fermions, the
  quantization of electric charge, and the absence of magnetic monopoles
  follow from admissibility constraints on the projection fiber.

  The framework yields falsifiable predictions: a~${\sim}5\%$ enhancement
  of~$H(z)$ at~$z \sim 1$ relative to~$\Lambda$CDM, environment-dependent
  deviations in kinematic inference in cosmic voids testable with current
  surveys, and a structural upper bound on the capacity exponent governing
  the fermion mass hierarchy.

  The present document serves as the conceptual and structural reference
  for the Cosmochrony programme.
  Formal derivations, spectral numerical evidence, and phenomenological
  applications are developed in a series of companion papers cited
  throughout.
}
